% Mathematical fragment for integration by the parent agent.
% Standard theorem environments, amsmath/amssymb, and \psi are assumed.
% Every label begins with pc:.

\section{Evaluated contacts for arbitrarily many polygamma factors}
\label{pc:section}

The latest incoming \emph{Polygamma Moments and Stieltjes Resonances}
report, its further-research question~6, asks for three independently shifted
Stieltjes factors \cite{PolygammaResonances}.  The following result answers the local-contact problem in the entire
Stieltjes-index-zero sector, with any number of factors and arbitrary
nonnegative argument derivative orders.  Its uniform comparison permits
all positive paths, including nonmonomial paths.  Spectral derivatives
corresponding to positive Stieltjes indices are not claimed here.
The finite-part integration and divided-difference tools are classical
\cite{DeBoor,TicaGalapon,VillanuevaGalapon};
the claimed contribution relative to the inspected sources is their
complete singular-subset assembly and the evaluated consequences below.

\subsection{The rational kernel as one Hermite divided difference}

For positive nodes $a_1,\ldots,a_p$, positive integers
$h_1,\ldots,h_p$, and $H=h_1+\cdots+h_p$, define, for
$0\le d\le H-1$,
\begin{equation}\label{pc:kernel-def}
 K_{\boldsymbol h,d}(\boldsymbol a)
 =\lim_{R\to\infty}\left\{
       \int_0^R\frac{x^d\,dx}{\prod_{i=1}^p(x+a_i)^{h_i}}
                 -\mathbf1_{d=H-1}\log R\right\}.
\end{equation}
Thus only the last numerator degree needs a finite part at infinity.
For a smooth function $F$ on the positive axis, write
$[a_1^{(h_1)},\ldots,a_p^{(h_p)}]F$ for its Hermite divided
difference with the node $a_i$ repeated $h_i$ times; its order is $H-1$.
The standard confluent convention \cite{DeBoor} gives
$[a^{(H)}]F=F^{(H-1)}(a)/(H-1)!$.

\begin{lemma}[Completely evaluated local kernel]\label{pc:kernel}
For every indicated degree and every positive set of nodes,
\begin{equation}\label{pc:kernel-dd}
 \boxed{\displaystyle
 K_{\boldsymbol h,d}(\boldsymbol a)
     =(-1)^{H+d}
       [a_1^{(h_1)},\ldots,a_p^{(h_p)}](z^d\log z).}
\end{equation}
In particular this is a finite expression in rational functions of the
nodes and their logarithms, including all confluent cases.
\end{lemma}

\begin{proof}
First take $H$ distinct nodes $b_1,\ldots,b_H$ and all multiplicities
equal to one.  The partial fraction coefficients of
$x^d/\prod_j(x+b_j)$ are
\[
 c_i=\frac{(-b_i)^d}{\prod_{j\ne i}(b_j-b_i)}.
\]
If $d\le H-2$ their sum is zero; if $d=H-1$ their sum is one.
Integrating from zero to $R$ and subtracting the prescribed logarithm
therefore gives
\[
 -\sum_i c_i\log b_i
   =(-1)^{H+d}\sum_i
          \frac{b_i^d\log b_i}{\prod_{j\ne i}(b_i-b_j)},
\]
which is exactly the divided difference.  Coalescing nodes proves the
general case.  For $d<H-1$ dominated convergence applies directly.
For $d=H-1$, apply it to the integral on $(0,1)$ and to
$x^d/\prod_j(x+b_j)-1/x$ on $(1,\infty)$, whose tail has a
locally uniform $O(x^{-2})$ bound.  This also proves the confluent
continuation without choosing singular individual coefficients.
\end{proof}

For direct implementation at distinct nodes one can instead use
\begin{align}\label{pc:partial-fractions}
 c_{i,q}&=[y^{h_i-q}]\,(y-a_i)^d
                      \prod_{j\ne i}(y+a_j-a_i)^{-h_j},\\
 K_{\boldsymbol h,d}(\boldsymbol a)
   &=-\sum_i c_{i,1}\log a_i
       +\sum_i\sum_{q=2}^{h_i}
                    \frac{c_{i,q}a_i^{1-q}}{q-1}.
                    \label{pc:kernel-rational}
\end{align}
Every coefficient in the first line is finite.  These equations are
also a transparent way to extract power--logarithm constants along
specified paths.

\subsection{The complete uniform local correction}

Fix $N\ge1$ and $r_1,\ldots,r_N\in\mathbb Z_{\ge0}$.  Put
\[
 h_i=r_i+1,\qquad p_i=(-1)^{r_i+1}r_i!,\qquad
 g_i(x)=\psi^{(r_i)}(1+x).
\]
The classical polygamma recurrence \cite[Eq.~5.15.5]{DLMF} gives
\begin{equation}\label{pc:split}
 \psi^{(r_i)}(x)=p_i x^{-h_i}+g_i(x).
\end{equation}
For a nonempty subset $S\subseteq\{1,\ldots,N\}$ let
\begin{align}
 H_S&=\sum_{i\in S}h_i,\qquad p_S=\prod_{i\in S}p_i,\\
 A_{S,d}(\boldsymbol a)
   &=[x^d]\prod_{j\notin S}g_j(x+a_j)
     =\sum_{\substack{d_j\ge0\ (j\notin S)\\\sum d_j=d}}
          \prod_{j\notin S}
              \frac{\psi^{(r_j+d_j)}(1+a_j)}{d_j!}.
              \label{pc:A}
\end{align}
The empty product convention makes $A_{S,0}=1$ and $A_{S,d}=0$
for $d>0$ when $S$ is the full set.  Crucially, the shifts in
\eqref{pc:A} are retained exactly.

Define the fixed-coordinate merged finite part
\begin{equation}\label{pc:merged}
 F_{\boldsymbol r}
    =\operatorname{CT}_{\delta\downarrow0}
                  \int_\delta^1\prod_{i=1}^N\psi^{(r_i)}(x)\,dx
\end{equation}
and the completely evaluated correction
\begin{align}\label{pc:Q}
 Q_{\boldsymbol r}(\boldsymbol a)
   &=\sum_{\varnothing\ne S\subseteq\{1,\ldots,N\}}
        p_S\sum_{d=0}^{H_S-1}A_{S,d}(\boldsymbol a)
                        K_{\boldsymbol h_S,d}(\boldsymbol a_S)\notag\\
   &=\sum_{\varnothing\ne S}
       \left(\prod_{i\in S}r_i!\right)
       \sum_{d=0}^{H_S-1}(-1)^d A_{S,d}(\boldsymbol a)
        [a_i^{(h_i)}:i\in S](z^d\log z).
\end{align}

\begin{theorem}[Every polygamma collision with positive shifts]
\label{pc:uniform}
As all $a_i\downarrow0$, with no restriction on their relative rates,
\begin{equation}\label{pc:uniform-formula}
 \boxed{\displaystyle
 \int_0^1\prod_{i=1}^N\psi^{(r_i)}(x+a_i)\,dx
          =F_{\boldsymbol r}+Q_{\boldsymbol r}(\boldsymbol a)+o(1).}
\end{equation}
The $o(1)$ is uniform as $\max_i a_i\to0$ through positive shifts.
The same statement holds after replacing any zero-order factor
$\psi$ by $\psi+c_i$, provided the same constant is added to its
regular part in \eqref{pc:A} and in the merged finite part.
\end{theorem}

\begin{proof}
Expand the product by choosing its singular part on a subset $S$.
For $S\ne\varnothing$, put
\[
 R_S(x;\boldsymbol a)=\prod_{i\in S}(x+a_i)^{-h_i},\qquad
 G_S(x;\boldsymbol a)=\prod_{j\notin S}g_j(x+a_j).
\]
Uniform Taylor expansion on a fixed compact interval gives
\[
 G_S(x;\boldsymbol a)
    =\sum_{d=0}^{H_S-1}A_{S,d}(\boldsymbol a)x^d
                   +x^{H_S}B_S(x;\boldsymbol a),
\]
where $B_S$ is uniformly bounded and tends to $B_S(x;0)$ for each
$x\in[0,1]$.  Since $R_S(x;\boldsymbol a)\le x^{-H_S}$ for
positive shifts, the integral of $x^{H_S}R_SB_S$ tends uniformly
to $\int_0^1B_S(x;0)dx$.  This estimate does not compare one shift
with another.

For $d\le H_S-2$, comparison with the tail on $(1,\infty)$ yields
\[
 \int_0^1x^dR_S(x;\boldsymbol a)dx
       =K_{\boldsymbol h_S,d}(\boldsymbol a_S)
                              -\frac1{H_S-d-1}+o(1).
\]
For $d=H_S-1$, compare instead with
$\int_1^\infty(x^{H_S-1}R_S-1/x)dx$; this tends to zero uniformly,
so the same formula holds with its rational term omitted.  Both
tail comparisons follow from an $O(\max a_i)x^{-2}$ bound after
the appropriate leading power is removed.

The empty-subset term converges ordinarily.  The limits of all
remainders and all displayed rational terms sum exactly to
$F_{\boldsymbol r}$: a merged monomial $x^{d-H_S}$ contributes
$-1/(H_S-d-1)$ for $d\le H_S-2$, while $x^{-1}$ contributes
zero to the fixed-coordinate constant on $(0,1)$.  This proves
\eqref{pc:uniform-formula}.  Uniform dominated convergence follows
by taking the supremum over the shrinking parameter cube in the
same integrable bound.  Constants added to the regular parts
do not change the proof.
\end{proof}

This theorem evaluates the entire local correction.  It does not assert
a new scalar evaluation of every merged product $F_{\boldsymbol r}$;
some such merged values remain spectral or harmonic constants in the
prior manuscript.  The next family eliminates that baseline by an
ordinary antiderivative and hence gives complete scalar identities.

\subsection{Nonresonant scales and their ordering}

Assume
\begin{equation}\label{pc:scales}
 a_i=\lambda_i\epsilon^{\alpha_i},\qquad
 \alpha_1>\alpha_2>\cdots>\alpha_N>0,\qquad\lambda_i>0,
\end{equation}
and that $\alpha_1,\ldots,\alpha_N$ are linearly independent over
$\mathbb Q$.  For a nonempty subset $S$, write $m(S)=\max S$;
its node is the largest shift within $S$.  Let $H_0=0$ and
$H_r=1+\cdots+1/r$.

\begin{theorem}[Evaluated nonresonant constant]\label{pc:nonresonant}
The ordinary shifted integral has a power--logarithm expansion through
its constant, and
\begin{equation}\label{pc:nonresonant-constant}
 \boxed{\displaystyle
 \operatorname{CT}_{\epsilon}
      \int_0^1\prod_i\psi^{(r_i)}(x+a_i)dx
  =F_{\boldsymbol r}
    -\sum_{\varnothing\ne S}p_S A_{S,H_S-1}(0)
            \bigl(\log\lambda_{m(S)}+H_{r_{m(S)}}\bigr).}
\end{equation}
The formula includes the same centering modification as
Theorem~\ref{pc:uniform}.  Its correction is a finite polynomial in
ordinary zeta values, Euler's constant and the displayed logarithms.
\end{theorem}

\begin{proof}
The rational formula \eqref{pc:kernel-rational} expands in convergent
geometric series in ratios of smaller shifts to larger shifts.
After each finite pole order is factored out, all remaining increments
of the exponent are positive differences $\alpha_i-\alpha_j$ or
positive regular-part exponents.  Therefore only finitely many terms
can occur below any fixed power.  Analytic expansion of the regular
coefficients has the same property.

For $d<H_S-1$, the rational coefficients of $K_{\boldsymbol h_S,d}$
are homogeneous of degree $d+1-H_S<0$, and the logarithms are
$\log\lambda_i+\alpha_i\log\epsilon$.  A zero exponent would,
by rational independence, require every integer exponent of every
shift variable to be zero.  Powers coming from regular-part shifts
occur only in variables outside $S$, so they must all be zero.
The remaining degree cannot be zero, a contradiction.  Thus no
noncritical degree contributes a constant.

For $d=H_S-1$, regular-part powers outside $S$ again cannot contribute.
Put $b=a_{m(S)}$ and $\rho_i=a_i/b$.  Scaling the defining integral
gives
\[
 K_{\boldsymbol h_S,H_S-1}(\boldsymbol a_S)
     =-\log b+
          K_{\boldsymbol h_S,H_S-1}
                         ((\rho_i)_{i\ne m(S)},1).
\]
Dominated convergence on $(0,1)$ and on the compensated tail gives
\begin{align*}
 \lim_{\rho_i\to0}
 K_{\boldsymbol h_S,H_S-1}((\rho_i),1)
 &=\lim_{R\to\infty}\left\{
   \int_0^R\frac{x^{h_{m(S)}-1}}{(1+x)^{h_{m(S)}}}dx-\log R
                                        \right\}\\
 &=-H_{h_{m(S)}-1}.
\end{align*}
For the first dominated-convergence step use the majorant
$x^{h_{m(S)}-1}/(1+x)^{h_{m(S)}}$ on $(0,1)$; the compensated
tail is uniformly $O(x^{-2})$.  The last evaluation follows from
$t=x/(1+x)$ and
$t^{h-1}/(1-t)=1/(1-t)-\sum_{j=0}^{h-2}t^j$.
Consequently the critical constant is
$-\log\lambda_{m(S)}-H_{r_{m(S)}}$.  Substitution into
\eqref{pc:Q} proves the theorem.
\end{proof}

Rational independence is a sufficient clean hypothesis, not an assertion
that every rationally dependent choice is resonant.  The exact kernel
formula remains valid for every positive path and permits finite
coefficient extraction in every specified monomial case.

\subsection{A family with complete zeta-valued answers}

Let
\[
 \Psi(x)=\psi(x)+\gamma,\qquad
 h(x)=\psi(1+x)+\gamma,\qquad
 f(x)=xh(x)=\sum_{k\ge2}(-1)^k\zeta(k)x^k.
\]
The series is the classical digamma Taylor expansion
\cite[Eq.~5.7.4]{DLMF}.
Thus $\Psi(x)=-1/x+h(x)$ and $\Psi(1)=h(0)=0$.
Take unit amplitudes in \eqref{pc:scales}.  For $1\le r\le N$,
define the full constant of the ordinary interval integral
\begin{equation}\label{pc:T-def}
 T_{N,r}=\operatorname{CT}_{\epsilon}
       \int_0^1\psi'(x+\epsilon^{\alpha_r})
               \prod_{j\ne r}\Psi(x+\epsilon^{\alpha_j})\,dx.
\end{equation}

\begin{theorem}[Every centered one-derivative moment]
\label{pc:all-powers}
For every $N\ge2$ and $1\le r\le N$,
\begin{equation}\label{pc:all-powers-formula}
 \boxed{\displaystyle
 T_{N,r}=-[x^N]\left\{
           \frac{(f(x)-1)^N}{N}
               +f(x)^{N-r}(f(x)-1)^{r-1}\right\}.}
\end{equation}
This is a finite polynomial in ordinary zeta values of total weight $N$.
It is independent of the particular nonresonant exponents and depends
only on the position of the derivative in their ordering.  Moreover
\begin{equation}\label{pc:sum-jump}
 \sum_{r=1}^N T_{N,r}=0,\qquad
 T_{N,r+1}-T_{N,r}
       =[x^N]f(x)^{N-r-1}(f(x)-1)^{r-1}.
\end{equation}
\end{theorem}

\begin{proof}
The merged antiderivative is $\Psi(x)^N/N$, so its finite part is
\[
 B_N=-\frac1N[x^N](f(x)-1)^N.
\]
In \eqref{pc:nonresonant-constant}, only a singular subset that
contains the derivative, has that derivative as its largest shift,
and contains any selection of its $r-1$ smaller shifts can contribute:
all other subsets have harmonic number $H_0=0$.
A subset of size $s$ has $H_S=s+1$, prefactor $p_S=(-1)^{s-1}$,
and regular coefficient $[x^s]h(x)^{N-s}$.  Its contact is therefore
$(-1)^s[x^s]h(x)^{N-s}$, with multiplicity $\binom{r-1}{s-1}$.
Their sum is
\begin{align*}
 C_{N,r}
 &=\sum_{s=1}^r(-1)^s\binom{r-1}{s-1}
                           [x^s]h(x)^{N-s}\\
 &=-[x^N]f(x)^{N-r}(f(x)-1)^{r-1}.
\end{align*}
Adding $B_N$ gives the theorem.  The adjacent difference is immediate.
For the sum use
$\sum_{r=1}^N f^{N-r}(f-1)^{r-1}=f^N-(f-1)^N$ and
$[x^N]f^N=0$, since $f$ starts in degree two.
\end{proof}

The first rows, simplified with $\zeta(2)^2=\tfrac52\zeta(4)$, are
\begin{align}
 (T_{2,1},T_{2,2})&=(0,0),\\
 (T_{3,1},T_{3,2},T_{3,3})&=(\zeta(3),0,-\zeta(3)),
                                                  \label{pc:triple-row}\\
 (T_{4,1},T_{4,2},T_{4,3},T_{4,4})
   &=\frac{\zeta(4)}4(-11,-1,5,7),
                                                  \label{pc:fourth-row}\\
 (T_{5,1},\ldots,T_{5,5})
   &=\bigl(\zeta(5)-4\zeta(2)\zeta(3),
            \zeta(5)-4\zeta(2)\zeta(3),\notag\\
   &\qquad\zeta(5)-2\zeta(2)\zeta(3),
            2\zeta(2)\zeta(3),
            8\zeta(2)\zeta(3)-3\zeta(5)\bigr).
                                                  \label{pc:fifth-row}
\end{align}
These are full asymptotic constants of the integrals, not merely their
differences from the merged finite part.

\subsection{The uncentered triple and an explicit resonant path}

For the uncentered triple the merged value is
\begin{equation}\label{pc:triple-baseline}
 \operatorname{CT}_{\delta\downarrow0}
       \int_\delta^1\psi'(x)\psi(x)^2\,dx
           =\zeta(3)-2\gamma\zeta(2).
\end{equation}
The three nonresonant contacts, with unit amplitudes, are respectively
$2\gamma\zeta(2)$,
$2\gamma\zeta(2)-\zeta(3)$, and
$2\gamma\zeta(2)-2\zeta(3)$.  Thus the full constants remain
exactly the triple row \eqref{pc:triple-row}.  A concrete admissible
choice is $(\alpha_1,\alpha_2,\alpha_3)=(\sqrt3,\sqrt2,1)$.
The cyclic sum is zero, also directly forced by integrating the
derivative of $\prod_i\psi(x+a_i)$ and taking its constant.

For comparison, the resonant path $(a,b,c)=(\epsilon^3,\epsilon^2,
\epsilon)$ with the derivative at $a$ has the different exact value
\begin{equation}\label{pc:integer-triple}
 \boxed{\displaystyle
 \operatorname{CT}_{\epsilon}
   \int_0^1\psi'(x+\epsilon^3)\psi(x+\epsilon^2)
                            \psi(x+\epsilon)dx
 =4+2\gamma-\zeta(2)+\zeta(3)
       +\left(\frac92-\gamma\right)\zeta(4)
       +\zeta(5)-\zeta(6).}
\end{equation}
Here is a finite certificate of its constant extraction.
At unit amplitudes every logarithm in \eqref{pc:kernel-rational}
is an integer multiple of $\log\epsilon$, so only its rational part
can contribute to the log-free constant.  Any singular subset not
containing the derivative has only simple poles and hence contributes
zero.  The four other subsets give the following constants, where
$h(x)=\psi(1+x)$:
\begin{center}
\begin{tabular}{c|l}
Singular subset & Its contact constant\\\hline
$\{a\}$ & $-\gamma\zeta(4)+\zeta(2)^2+2\gamma\zeta(2)$\\
$\{a,b\}$ & $\gamma-\zeta(3)+2\zeta(4)+\zeta(5)-\zeta(6)$\\
$\{a,c\}$ & $\gamma-\zeta(2)+\zeta(3)$\\
$\{a,b,c\}$ & $4$
\end{tabular}
\end{center}
For example the last rational part is
$1/[a(b-a)(c-a)]=\epsilon^{-6}/[(1-\epsilon)(1-\epsilon^2)]$,
whose constant is four.  The two-factor subset with remaining
regular factor $h(x+c)$ has rational part
\[
 -\frac{h(c)-a h'(c)+a^2h''(c)/2}{a(b-a)},
\]
and the exchanged expression gives the other row.  The singleton
has rational part
$h(b)h(c)/a-h'(b)h(c)-h(b)h'(c)$.
Taking the indicated finite Taylor coefficients and adding
\eqref{pc:triple-baseline} proves \eqref{pc:integer-triple}.

\paragraph{Ordinary renormalized limits.}
The symbol $\operatorname{CT}_{\epsilon}$ always selects
$\epsilon^0(\log\epsilon)^0$ from a proved expansion; it is not the
ordinary limit of a divergent integral.  All the results also give
ordinary limits: subtract the finite expression
$Q_{\boldsymbol r}(\boldsymbol a)$ from the original integral and
then add its evaluated constant.  Equation~\eqref{pc:uniform-formula}
proves the resulting limit equals the stated full constant.  Thus the
subtraction is an explicit finite formula in polygammas, rational
functions and logarithms, with no fitted coefficients or new integrals.


The exact and numerical verification records for this section are
described in Section~\ref{repro:section}. The kernel methods are
classical \cite{DeBoor,Lesch,TicaGalapon,VillanuevaGalapon};
the complete assembly and evaluated families are the continuation
results relative to the inspected project sources.
